\documentclass[10pt,a4paper]{article}
\usepackage[a4paper,margin=14mm,top=12mm,bottom=14mm]{geometry}
\usepackage{fontspec}
\setmainfont{Liberation Sans}
\usepackage{polyglossia}
\setdefaultlanguage{french}
\usepackage{amsmath}
\usepackage{xcolor}
\usepackage{graphicx}
\usepackage[most]{tcolorbox}
\usepackage{tabularx,booktabs,array,multirow}
\usepackage{enumitem}
\usepackage{fancyhdr}
\usepackage{lastpage}
\usepackage[hidelinks]{hyperref}
\definecolor{uimm}{HTML}{005677}
\definecolor{cyan}{HTML}{13A2DD}
\definecolor{fond}{HTML}{E6F2FA}
\definecolor{ok}{HTML}{2E7D32}
\definecolor{ko}{HTML}{B3261E}
\definecolor{sylvain}{HTML}{7B1FA2}
\setlength{\parindent}{0pt}
\setlength{\parskip}{3pt}
\setlist{nosep}
\setlist[itemize]{leftmargin=13pt,label=\textcolor{cyan}{\textbullet}}
\newcommand{\titre}[1]{\vspace{5pt}{\color{uimm}\large\bfseries #1}\par\vspace{-3pt}{\color{cyan}\rule{\linewidth}{1.2pt}}\par\vspace{0pt}}
\tcbset{boite/.style={colback=fond,colframe=uimm,boxrule=0.6pt,arc=2pt,left=6pt,right=6pt,top=4pt,bottom=4pt},
        defi/.style={colback=sylvain!6,colframe=sylvain,boxrule=0.8pt,arc=2pt,left=6pt,right=6pt,top=4pt,bottom=4pt,
                     title={\bfseries Défi},coltitle=white,colbacktitle=sylvain,fonttitle=\small}}
\newcommand{\bandeau}[3]{%
\begin{tcolorbox}[colback=uimm,colframe=uimm,arc=2pt,left=4pt,right=8pt,top=4pt,bottom=4pt]
\begin{minipage}[c]{18mm}\setlength{\fboxsep}{1.2pt}\fcolorbox{white}{white}{\includegraphics[height=16.5mm]{logo_uimm.jpg}}\end{minipage}\hspace{4mm}%
\begin{minipage}[c]{\dimexpr\linewidth-22mm\relax}\color{white}
{\small #1}\\[1pt]{\LARGE\bfseries #2}\\[2pt]{\large #3}\hfill{\small CQPM Technicien qualité}
\end{minipage}
\end{tcolorbox}}
\newcommand{\reponse}[1][1]{\par\vspace{2pt}{\color{gray!60}\foreach \i in {1,...,#1}{\rule{\linewidth}{0.3pt}\par\vspace{9pt}}}}
\usepackage{pgffor}
\pagestyle{fancy}\fancyhf{}
\renewcommand{\headrulewidth}{0pt}
\newcommand{\pied}[1]{\fancyfoot[L]{\footnotesize\color{uimm}#1}\fancyfoot[R]{\footnotesize\color{uimm}Page \thepage{} sur \pageref{LastPage}}}

\pied{TP1 · L'instrument est-il capable ? · fiche stagiaire}
\begin{document}
\thispagestyle{fancy}
\bandeau{Travaux pratiques · Capabilité du moyen de mesure}{TP1 · L'instrument est-il capable ?}{Étude de type 1 sur cale étalon, méthode Bosch}

\begin{tcolorbox}[boite]
\textbf{La question.} Pour contrôler le diamètre de la butée (tolérance $T = 0{,}15$ mm), peut-on se fier au \textbf{micromètre} ? Et au \textbf{pied à coulisse} ?
On mesure plusieurs fois une cale étalon dont on connaît la vraie valeur, puis on calcule Cg et Cgk.
\end{tcolorbox}

\titre{1. Le rappel}
\begin{minipage}[c]{0.52\linewidth}
\[C_g=\frac{0{,}2\,T}{6\,s}\qquad C_{gk}=\frac{0{,}1\,T-\lvert\bar x-x_m\rvert}{3\,s}\]
\end{minipage}\hfill
\begin{minipage}[c]{0.45\linewidth}\small
$\bar x$ : moyenne des mesures ; $s$ : écart-type ;\\ $x_m$ : valeur certifiée de la cale.\\
\textbf{Capable si $C_g \ge 1{,}33$ et $C_{gk} \ge 1{,}33$.}
\end{minipage}

{\small Avant tout calcul : la résolution de l'instrument doit rester sous 5 \% de $T$. Voir la fiche « Le moyen de mesure est-il capable ? ».}

\titre{2. Le matériel}
\renewcommand{\arraystretch}{1.35}
\begin{tabularx}{\linewidth}{@{}l X@{}}
Cale étalon & Valeur certifiée $x_m$ = \rule{25mm}{0.3pt} mm \quad Certificat n° \rule{25mm}{0.3pt}\\
Micromètre & Résolution : 0,001 mm \quad N° \rule{25mm}{0.3pt}\\
Pied à coulisse & Résolution : 0,02 mm (au 1/50) \quad N° \rule{25mm}{0.3pt}\\
Calculs & Page web « L'instrument est-il capable ? »\\
\end{tabularx}

\titre{3. Le protocole}
\begin{enumerate}[label=\textcolor{cyan}{\textbf{\arabic*.}},leftmargin=15pt]
\item En binôme : l'un mesure, l'autre note. On échange les rôles en changeant d'instrument.
\item Nettoyer la cale. La tenir par les bords, sans la serrer dans la main : la chaleur des doigts la dilate.
\item \textbf{Partie A : micromètre.} Mesurer la cale \textbf{25 fois} (le livret Bosch demande 50 mesures, au moins 25). Utiliser le cliquet. \textbf{Retirer puis remettre} la cale avant chaque mesure.
\item \textbf{Partie B : pied à coulisse.} Même chose, 25 mesures.
\item Noter toutes les valeurs, \textbf{n'en supprimer aucune}, même « bizarre ».
\item Saisir chaque série dans la page, avec $T = 0{,}15$, la valeur certifiée et la résolution de l'instrument.
\end{enumerate}

\titre{4. Les relevés (mm)}
\renewcommand{\arraystretch}{1.6}
\textbf{Partie A : micromètre}\hfill Mesuré par \rule{30mm}{0.3pt}\\[2pt]
\begin{tabularx}{\linewidth}{|*{10}{>{\centering\arraybackslash}X|}}\hline
\scriptsize 1 & & \scriptsize 2 & & \scriptsize 3 & & \scriptsize 4 & & \scriptsize 5 & \\\hline
\scriptsize 6 & & \scriptsize 7 & & \scriptsize 8 & & \scriptsize 9 & & \scriptsize 10 & \\\hline
\scriptsize 11 & & \scriptsize 12 & & \scriptsize 13 & & \scriptsize 14 & & \scriptsize 15 & \\\hline
\scriptsize 16 & & \scriptsize 17 & & \scriptsize 18 & & \scriptsize 19 & & \scriptsize 20 & \\\hline
\scriptsize 21 & & \scriptsize 22 & & \scriptsize 23 & & \scriptsize 24 & & \scriptsize 25 & \\\hline
\end{tabularx}

\vspace{6pt}
\textbf{Partie B : pied à coulisse}\hfill Mesuré par \rule{30mm}{0.3pt}\\[2pt]
\begin{tabularx}{\linewidth}{|*{10}{>{\centering\arraybackslash}X|}}\hline
\scriptsize 1 & & \scriptsize 2 & & \scriptsize 3 & & \scriptsize 4 & & \scriptsize 5 & \\\hline
\scriptsize 6 & & \scriptsize 7 & & \scriptsize 8 & & \scriptsize 9 & & \scriptsize 10 & \\\hline
\scriptsize 11 & & \scriptsize 12 & & \scriptsize 13 & & \scriptsize 14 & & \scriptsize 15 & \\\hline
\scriptsize 16 & & \scriptsize 17 & & \scriptsize 18 & & \scriptsize 19 & & \scriptsize 20 & \\\hline
\scriptsize 21 & & \scriptsize 22 & & \scriptsize 23 & & \scriptsize 24 & & \scriptsize 25 & \\\hline
\end{tabularx}

\newpage
\titre{5. Les résultats}
\renewcommand{\arraystretch}{1.6}
\begin{tabularx}{\linewidth}{|l|*{6}{>{\centering\arraybackslash}X|}}\hline
 & \textbf{Résolution en \% de T} & \textbf{$\bar x$} & \textbf{$s$} & \textbf{Biais $\bar x - x_m$} & \textbf{$C_g$} & \textbf{$C_{gk}$}\\\hline
Micromètre & & & & & & \\\hline
Pied à coulisse & & & & & & \\\hline
\end{tabularx}

\titre{6. Les questions}
\begin{enumerate}[label=\textcolor{cyan}{\textbf{Q\arabic*.}},leftmargin=22pt]
\item Calculez la résolution de chaque instrument en pourcentage de $T$. Lequel passe le critère des 5 \% ?\reponse
\item Le micromètre est-il capable pour la butée ? Justifiez avec $C_g$ et $C_{gk}$.\reponse
\item Le pied à coulisse est-il capable ? Que penser de son résultat si toutes ses mesures sont identiques (que vaut alors $s$) ?\reponse[2]
\item Pour l'un des instruments, $C_g \ge 1{,}33$ mais $C_{gk} < 1{,}33$. Que se passe-t-il ? Que faut-il vérifier ?\reponse[2]
\item Quel instrument choisir pour contrôler la butée au poste ? Pourquoi ?\reponse
\end{enumerate}

\begin{tcolorbox}[defi]
\begin{enumerate}[label=\textbf{\alph*.},leftmargin=14pt]
\item \textbf{Le biais est-il réel ?} Le livret (annexe C.1) considère le biais comme négligeable, au niveau de confiance de 95 \%, si
\[\frac{\lvert\bar x-x_m\rvert}{s}\le 0{,}413\ \text{pour 25 mesures}\qquad(0{,}284\ \text{pour 50 mesures}).\]
Appliquez ce test au micromètre. Montrez qu'il revient à vérifier $C_g - C_{gk} \le 0{,}138$ pour 25 mesures.
\item \textbf{La plus petite tolérance contrôlable.} À partir des formules de $C_g$ et $C_{gk}$, exprimez la plus petite tolérance $T_{\min}$ pour laquelle le micromètre reste capable, en fonction de $s$ et du biais. Calculez-la avec vos mesures.
\item \textbf{Aux deux bouts de la tolérance.} Si des cales le permettent, refaites la partie A avec une cale proche de 11,85 mm puis de 12,00 mm. Le biais est-il le même aux deux bouts ? (Le livret, § 4.4, recommande ces deux points quand on ne dispose que de deux étalons.)
\end{enumerate}
\end{tcolorbox}

\vfill
{\footnotesize\color{uimm}\textbf{Source.} Robert Bosch GmbH, \emph{Booklet 10 : Capability of Measurement and Test Processes}, édition novembre 2019 : § 4.1 (procédure 1), fig. 3 (résolution), § 4.4 (linéarité) et annexe C.1 (signification du biais).}
\end{document}
